\documentclass[11pt]{article}
\usepackage[margin=1.1in]{geometry}
\usepackage{amsmath,amssymb,amsthm}
\usepackage[colorlinks=true,linkcolor=blue,citecolor=blue,urlcolor=blue]{hyperref}

\newtheorem{theorem}{Theorem}[section]
\newtheorem{lemma}[theorem]{Lemma}
\newtheorem{proposition}[theorem]{Proposition}
\theoremstyle{remark}
\newtheorem{remarkc}[theorem]{Remark}

\newcommand{\Dfour}{\ensuremath{\mathsf{D4}}}
\newcommand{\UNK}{\ensuremath{\mathsf{UNK}}}
\newcommand{\FAL}{\ensuremath{\mathsf{FAL}}}
\newcommand{\TRU}{\ensuremath{\mathsf{TRU}}}
\newcommand{\CON}{\ensuremath{\mathsf{CON}}}
\newcommand{\meetk}{\ensuremath{\wedge_k}}
\newcommand{\joink}{\ensuremath{\vee_k}}
\newcommand{\Clo}{\ensuremath{\mathrm{Clo}}}
\newcommand{\GR}{\ensuremath{G|_R}}
\newcommand{\Dsame}{\ensuremath{D_{\mathrm{same}}}}
\newcommand{\Dopp}{\ensuremath{D_{\mathrm{opp}}}}

\title{The Resolved-Face Crown of the Dual-Rail Carrier:\\
Self-Dual Completeness, the Exact Count,\\ and One-Constant Completion}
\author{Won Chul Yang\\ \small Independent Researcher, Seoul, Republic of Korea\\
\small \texttt{wcy0969@gmail.com}}
\date{Public edition v1.0 --- 2026}

\begin{document}
\maketitle

\begin{center}\small\itshape
Companion note to ``The Price of NOT on \textsf{D4}'' (Paper~II,
DOI: 10.5281/zenodo.21800033): it proves in full the two crown facts that
Paper~II's Appendix~A.4--A.5 uses as boundary context and attributes to a
companion line. Every headline result is machine-verified in
dependency-minimal Lean~4 (core only; no mathlib; no
\texttt{native\_decide}; no \texttt{sorry}; kernel axiom profile at most
\texttt{[propext, Quot.sound]}).
\end{center}

\begin{abstract}
On the four-state dual-rail carrier \Dfour{} with closed core
$G=\Clo\{\meetk,\joink,P\}$, consider the resolved face $R=\{\FAL,\TRU\}$
and the outer fragment $\GR$ of $R$-valued restrictions of closed-core
terms on $R^n$. We prove the crown theorem: $\GR^{(n)}$ is exactly the
set of $P$-equivariant maps $R^n\to R$ --- equivalently, under the Boolean
reading of the face, exactly the self-dual Boolean functions of arity
$n$. The count is therefore $2^{2^{n-1}}$: the diagonal $P$-action on
$R^n$ is free, and an equivariant map freely chooses one output per
orbit. The self-dual reading corrects a natural miscount: $2^n$ counts
per-coordinate routing data, not functions --- odd-arity majority is
self-dual and lies in the crown, while AND, OR, and both constants do
not. We then prove one-constant completion: adjoining a single resolved
constant to the crown realizes every map $R^n\to R$, by an explicit
self-dualization with one guard variable. Finally we record why the crown
is a genuinely dual-rail (single-orbit) phenomenon: for resolved faces
with two or more $P$-orbits, $P$-equivariance strictly over-approximates
the outer fragment, witnessed by an explicit unary counterexample.
\end{abstract}

\noindent\textbf{Keywords:} clone theory; self-dual Boolean functions;
dual-rail carriers; equivariant maps; functional completeness; Lean~4.\\
\textbf{MSC 2020:} 03B50, 06D30, 08A40, 94C10.

\section{Introduction}

The companion paper \cite{wcy2} studies the four-state dual-rail carrier
$\Dfour=\{\UNK,\FAL,\TRU,\CON\}$ with closed core
$G=\Clo\{\meetk,\joink,P\}$ and proves the exact negation-complexity law
$\nu(f)=\mathrm{dec}(f)$ on the resolved face $R=\{\FAL,\TRU\}$ after the
confinement-breaking operation $\sigma_s$ is adjoined. Before $\sigma_s$
enters, the closed core already induces an \emph{outer fragment} on the
resolved face:
\[
\GR^{(n)} \;=\; \{\, t|_{R^n} \;:\; t \text{ an $n$-ary term of } G,\
t(R^n)\subseteq R \,\},
\]
the family of $R$-valued restrictions of closed-core terms. Paper~II's
Appendix~A.4--A.5 uses two facts about this fragment as boundary context
--- that it consists exactly of the $P$-equivariant maps and has
cardinality $2^{2^{n-1}}$, and that one adjoined resolved constant
completes it to all maps $R^n\to R$ --- attributing both to a companion
line. This note is that companion: it states and proves both facts in
full, together with the structural reading that makes them stable, and it
machine-verifies the results in Lean~4 on top of the published
mechanization of \cite{wcy2}.

Three points deserve emphasis.

\paragraph{The self-dual reading.} Under the Boolean encoding $\TRU=1$,
$\FAL=0$ of the face, $P$ acts as negation, and $P$-equivariance of
$h:R^n\to R$ becomes $h(\neg x)=\neg h(x)$: the crown is exactly the
clone-theoretic class of \emph{self-dual} Boolean functions
\cite{post,lau}. This immediately calibrates what the crown does and does
not exclude: AND, OR, and the constants are not self-dual and are
excluded; odd-arity majority \emph{is} self-dual and is included,
realized inside the carrier as orbit-dependent routing
(Proposition~\ref{prop:calib} gives the explicit ternary-majority term).
In particular, the correct count of the crown is the count of self-dual
functions, $2^{2^{n-1}}$ --- not $2^n$, which counts per-coordinate
routing data rather than functions; at $n=3$ the crown has 16 elements,
including majority, against $8=2^3$ (Remark~\ref{rem:miscount}).

\paragraph{The two-layer proof.} Necessity is inherited: closed-core
terms are $P$-equivariant on all of $\Dfour^n$ by the qualitative
soundness theorem of \cite{wcy2}, so their $R$-valued restrictions are
$P$-equivariant on the face. Sufficiency is an explicit construction:
one orbit-literal minterm per $\TRU$-point of the target map, joined; on
the face, the escape values \UNK{} of the mismatched minterms are
absorbed by the joins, and the single matched minterm (or the matched
anti-minterm) delivers the resolved output. No step uses $\sigma_s$.

\paragraph{The boundary.} The crown is a single-orbit phenomenon. On
resolved faces with two or more $P$-orbits, a resolved value carries both
a within-pair bit and an orbit label; $P$-equivariance constrains only
the former, while closed-core terms are confined to orbits.
Section~\ref{sec:sharp} records the minimal unary counterexample. This is
why the crown theorem is stated for the dual-rail carrier and claimed for
nothing larger.

\section{Preliminaries}

We use the setting of \cite[\S2]{wcy2}. $\Dfour$ is identified with
$\{0,1\}^2$ via $\UNK=(0,0)$, $\FAL=(0,1)$, $\TRU=(1,0)$, $\CON=(1,1)$;
$\meetk$ and $\joink$ are componentwise meet and join in the knowledge
order; $P(a,b)=(b,a)$. The resolved face is $R=\{\FAL,\TRU\}$, on which
$P$ restricts to the transposition. Under the Boolean encoding $\TRU=1$,
$\FAL=0$, a map $h:R^n\to R$ is $P$-equivariant exactly when it is
self-dual: $h(\neg x)=\neg h(x)$.

We use one published fact from \cite{wcy2}, machine-verified there:

\begin{itemize}
\item[\textbf{(E)}] (Theorem~2.1 of \cite{wcy2}) every term operation of
$G$ is monotone for the knowledge order, preserves both endpoints, and is
$P$-equivariant on $\Dfour^n$.
\end{itemize}

We also use the elementary evaluation table of the internal operations on
resolved and endpoint values (immediate from the componentwise
definitions, and machine-checked in the present artifact):
\[
\TRU\meetk\TRU=\TRU,\quad \FAL\meetk\FAL=\FAL,\quad \TRU\meetk\FAL=\UNK,
\]
\[
\UNK\meetk v=\UNK,\quad \TRU\joink\UNK=\TRU,\quad \FAL\joink\UNK=\FAL,
\quad \UNK\joink\UNK=\UNK,
\]
together with $P$-equivariance of $\meetk,\joink$ and $P(\FAL)=\TRU$.

\section{The Crown Theorem}

\begin{theorem}[crown; self-dual completeness of the outer fragment]
\label{thm:crown}
For every arity $n\ge 1$,
\[
\GR^{(n)} = \{\, h:R^n\to R \mid h \text{ is } P\text{-equivariant} \,\}
= \{\, \text{self-dual Boolean functions of arity } n \,\}.
\]
\end{theorem}

\begin{proof}
The second equality is the Boolean reading of Section~2. We prove the
first.

($\subseteq$) Let $t$ be an $n$-ary closed-core term with
$t(R^n)\subseteq R$. By (E), $t$ is $P$-equivariant on $\Dfour^n$;
restricting to $R^n$ (which is $P$-invariant) gives a $P$-equivariant map
$R^n\to R$.

($\supseteq$) Let $h:R^n\to R$ be $P$-equivariant. First, $h$ attains
the value \TRU: $P$-equivariance gives $h(P\vec x)=P\,h(\vec x)\ne
h(\vec x)$, so $h$ is nonconstant, and a map into $R$ that is nonconstant
attains both values. Let $A=\{\vec y\in R^n : h(\vec y)=\TRU\}$; $A$ is
nonempty, and $P(A)$ is exactly the set where $h=\FAL$ (equivariance), so
$A$ contains no $P$-pair: $\vec y\in A \Rightarrow P\vec y\notin A$.

For each $\vec y\in A$ define the \emph{orbit-literal minterm}
\[
M_{\vec y}(\vec x) \;=\; \bigwedge\nolimits_k \{\ell_1(x_1),\dots,
\ell_n(x_n)\},
\qquad
\ell_i(x_i)=
\begin{cases}
x_i, & y_i=\TRU,\\
P x_i, & y_i=\FAL,
\end{cases}
\]
a closed-core term (only $\meetk$ and $P$ occur). Evaluate $M_{\vec y}$
on a face point $\vec z\in R^n$. Each literal value $\ell_i(z_i)$ is
resolved, and equals \TRU{} exactly when $z_i=y_i$. On resolved values,
$\TRU\meetk\TRU=\TRU$, $\FAL\meetk\FAL=\FAL$, $\TRU\meetk\FAL=\UNK$, and
\UNK{} absorbs. Hence:
\begin{itemize}
\item if $\vec z=\vec y$, every literal is \TRU{} and
  $M_{\vec y}(\vec z)=\TRU$;
\item if $\vec z=P\vec y$, every literal is \FAL{} and
  $M_{\vec y}(\vec z)=\FAL$;
\item otherwise the literals are mixed and $M_{\vec y}(\vec z)=\UNK$.
\end{itemize}

Now set $t=\bigvee_k\{M_{\vec y} : \vec y\in A\}$ (a join over the
nonempty family $A$; again a closed-core term). Evaluate $t$ on
$\vec z\in R^n$:
\begin{itemize}
\item if $h(\vec z)=\TRU$, then $\vec z\in A$, so the minterm
  $M_{\vec z}$ contributes \TRU; every other $\vec y\in A$ has
  $\vec y\ne\vec z$ and also $P\vec y\ne\vec z$ (else $\vec z=P\vec y$
  with $\vec y\in A$, so $h(\vec z)=h(P\vec y)=P\,h(\vec y)=\FAL$,
  contradiction), so every other minterm contributes \UNK. Since
  $\TRU\joink\UNK=\TRU$ and $\UNK\joink\UNK=\UNK$, the join is
  $\TRU=h(\vec z)$.
\item if $h(\vec z)=\FAL$, then $P\vec z\in A$, so the minterm
  $M_{P\vec z}$ evaluates at $\vec z=P(P\vec z)$ to \FAL; every other
  $\vec y\in A$ has $\vec y\ne\vec z$ (as $h(\vec y)=\TRU$) and
  $\vec y\ne P\vec z$, so every other minterm contributes \UNK. Since
  $\FAL\joink\UNK=\FAL$ and no minterm contributes \TRU{} (that would
  require $\vec z\in A$), the join is $\FAL=h(\vec z)$.
\end{itemize}
In both cases $t(\vec z)=h(\vec z)\in R$. Hence $t|_{R^n}=h$ with
$t(R^n)\subseteq R$, so $h\in\GR^{(n)}$.
\end{proof}

\begin{remarkc}
The construction realizes $h$ with a term whose only generators are
$\meetk,\joink,P$ --- one orbit-literal minterm per \TRU-point of $h$. An
alternative route extends the face labeling monotonically through the
ambient rail cube (the face is a one-hot antichain in $\{0,1\}^{2n}$) and
lifts the resulting positive formula; the direct construction above is
the one mechanized.
\end{remarkc}

\begin{theorem}[the exact count]\label{thm:count}
For every $n\ge 1$, $\;|\GR^{(n)}| = 2^{2^{n-1}}$.
\end{theorem}

\begin{proof}
By Theorem~\ref{thm:crown} it suffices to count $P$-equivariant maps
$h:R^n\to R$. The diagonal $P$-action on $R^n$ is free: $P$ has no fixed
point on $R$, so $P\vec x=\vec x$ is impossible already in the first
coordinate. Hence $R^n$ (of size $2^n$) splits into $2^{n-1}$ orbits of
size two; each orbit has a unique representative with first coordinate
\TRU. A $P$-equivariant $h$ is freely determined by its values on the
$2^{n-1}$ representatives --- the value on the partner point is forced to
be the $P$-image --- and every such assignment defines an equivariant
map. This is a bijection between $\GR^{(n)}$ and the maps from the
representatives to $R$, of cardinality $2^{2^{n-1}}$.
\end{proof}

For $n=1,\dots,5$ the counts are $2,\,4,\,16,\,256,\,65{,}536$.

\section{What the Crown Contains --- and the Corrected Count}

\begin{proposition}[calibration]\label{prop:calib}
Under the Boolean reading:
\begin{enumerate}
\item AND, OR, and both constants do not lie in \GR{} (none is
  self-dual; constants are moreover not $P$-equivariant).
\item Odd-arity majority lies in \GR. Explicitly, for $n=3$, with the
  agreement detectors
  \[
  \Dsame(x,y)=(x\meetk y)\joink(Px\meetk Py),
  \qquad
  \Dopp(x,y)=\Dsame(x,Py),
  \]
  which are $E$-valued on $R^2$ ($E=\{\UNK,\CON\}$) with
  $\Dsame=\CON$ iff $x=y$ and $\Dopp=\CON$ iff $x=Py$, the term
  \[
  \mathrm{MAJ}_3(x_1,x_2,x_3)=
  \bigl(x_1\meetk(\Dsame(x_1,x_2)\joink\Dsame(x_1,x_3))\bigr)
  \joink
  \bigl(Px_1\meetk(\Dopp(x_1,x_2)\meetk\Dopp(x_1,x_3))\bigr)
  \]
  is a closed-core term whose restriction to $R^3$ is ternary majority.
\end{enumerate}
\end{proposition}

\begin{proof}
(1) Self-duality $f(\neg x)=\neg f(x)$ fails for AND at $x=(1,0)$ (it
would force $\mathrm{AND}(0,1)=\neg\,\mathrm{AND}(1,0)$, i.e.\
$0=\neg 0$), and dually for OR; constants $c$ satisfy $Pc\ne c$ on $R$,
so $P$-equivariance fails at once. (2) Majority of odd arity satisfies
$\mathrm{MAJ}(\neg x_1,\dots,\neg x_m)=\neg\,\mathrm{MAJ}(x_1,\dots,x_m)$
(complementing all inputs complements the strict majority), so it is
self-dual and lies in the crown by Theorem~\ref{thm:crown}. For the
explicit term: on resolved inputs the detectors are $E$-valued as stated,
by direct evaluation of the four resolved pairs (Proposition~A.1 of
\cite{wcy2}). On $E$, $\meetk$ and $\joink$ act as Boolean AND/OR with
$\CON=$ true, $\UNK=$ false. The first branch is active (its guard is
\CON) exactly when $x_2=x_1$ or $x_3=x_1$, in which case the majority
value is $x_1$; the second branch is active exactly when both
$x_2\ne x_1$ and $x_3\ne x_1$, in which case the majority value is
$Px_1$. The guard values absorb correctly on $R$:
$\UNK\meetk r=\UNK$, $\CON\meetk r=r$, $r\joink\UNK=r$, so the inactive
branch contributes \UNK{} and the active branch passes its routed value
through. Exhaustive evaluation over $R^3$ (machine-checked;
Section~\ref{sec:mech}) confirms the restriction is majority.
\end{proof}

\begin{remarkc}[the corrected count: routing data versus functions]
\label{rem:miscount}
A natural first count of the crown reads each coordinate as carrying an
independent forward/backward routing choice and arrives at $2^n$. That
figure counts \emph{routing data}, not functions: already at $n=3$ the
crown contains 16 functions (Theorem~\ref{thm:count}) --- including
ternary majority (Proposition~\ref{prop:calib}), which is not of the form
$\vec x\mapsto x_i$ or $\vec x\mapsto Px_i$ for any fixed $i$ ---
strictly more than $2^3=8$. An earlier internal manuscript in the
author's series stated the $2^n$ figure for this fragment; the present
note supersedes it. The stable statement is Theorem~\ref{thm:crown}: the
crown is self-dual completeness, and its size is the size of the
self-dual class.
\end{remarkc}

\section{One-Constant Completion}

\begin{theorem}[one resolved constant completes the outer fragment]
\label{thm:completion}
Fix either resolved constant, say $c=\FAL$. For every map $g:R^n\to R$
(not necessarily equivariant) there exist a $P$-equivariant map
$\hat g:R^{n+1}\to R$ with
\[
g(\vec x) = \hat g(c,\vec x),
\]
where $\hat g\in\GR^{(n+1)}$ by Theorem~\ref{thm:crown}. Hence the clone
generated on the face by \GR{} together with the constant $c$ contains
all $2^{2^n}$ maps $R^n\to R$.
\end{theorem}

\begin{proof}
In the Boolean reading ($c=0$), define the self-dualization of $g$ with
one guard variable:
\[
\hat g(x_0,\vec x)=
\begin{cases}
g(\vec x), & x_0=0,\\
\neg g(\neg\vec x), & x_0=1.
\end{cases}
\]
$\hat g$ is self-dual: complementing all $n+1$ arguments swaps the two
branches and complements the output ---
\[
\hat g(\neg x_0,\neg\vec x)=
\begin{cases}
g(\neg\vec x), & x_0=1,\\
\neg g(\vec x), & x_0=0,
\end{cases}
\]
which in both cases equals $\neg\hat g(x_0,\vec x)$. Specializing the
guard to the constant $0$ gives $\hat g(0,\vec x)=g(\vec x)$. By
Theorem~\ref{thm:crown}, $\hat g$ is the restriction of a closed-core
term $t$ of arity $n+1$; substituting the constant $c$ for the guard slot
of $t$ (a composition in the clone generated by $\GR\cup\{c\}$) realizes
$g$ on $R^n$. The count $2^{2^n}$ is the number of all maps $R^n\to R$.
\end{proof}

\begin{remarkc}
Theorem~\ref{thm:completion} is the zero/one accessibility contrast used
by \cite[App.~A.5]{wcy2}: without constants the outer fragment is the
proper self-dual class; a single resolved constant jumps it to
everything. (In clone-theoretic terms this reflects the maximality of the
self-dual clone in Post's lattice \cite{post}; the proof above is
self-contained and does not invoke the classification.)
\end{remarkc}

\section{Sharpness: the Crown Is a Single-Orbit Phenomenon}
\label{sec:sharp}

\begin{remarkc}[failure for larger resolved faces]
Suppose the resolved face has $2k\ge 4$ elements,
$R=O_1\sqcup\dots\sqcup O_k$ with $P$-pairs $O_i=\{a_i,b_i\}$. A resolved
value now carries two pieces of information: the within-pair bit and the
orbit label. $P$-equivariance constrains only the former. Closed-core
terms, however, are orbit-confined: a unary term evaluated at a resolved
$r$ can output only values in $\{r,Pr\}$ together with the endpoints,
because the generators $\meetk,\joink,P$ cannot manufacture an absolute
orbit label from a single input. Concretely, on the six-element flat
carrier with $R=\{a,b,c,d\}$, $P(a)=b$, $P(c)=d$, the unary map
\[
h(a)=a,\quad h(b)=b,\quad h(c)=a,\quad h(d)=b
\]
(extended by fixing the endpoints) is monotone, endpoint-preserving, and
$P$-equivariant, yet no closed-core term restricts to it: unary
closed-core terms act in one of the four global modes $x$, $Px$,
$x\meetk Px$, $x\joink Px$, and none maps $c$ into $\{a,b\}$. Hence for
$|R|\ge 4$, $P$-equivariance strictly over-approximates the outer
fragment, and the correct description is a routing/signature theory
rather than the crown. The crown theorem is therefore stated here only
for the dual-rail carrier, where $R$ is a single $P$-orbit and the two
notions coincide.
\end{remarkc}

\section{Mechanization}\label{sec:mech}

The note ships with a Lean~4 development that extends the published
artifact of \cite{wcy2} (core Lean~4 only; no mathlib; no
\texttt{native\_decide}; no \texttt{sorry}; toolchain pinned; kernel
axiom profile at most \texttt{[propext, Quot.sound]}):
\begin{itemize}
\item \textbf{necessity} --- \texttt{crown\_necessity}, from the
  mechanized Theorem~2.1 of \cite{wcy2} (\texttt{cterm\_equivariant});
\item \textbf{sufficiency} --- \texttt{crown\_sufficiency}: the
  orbit-literal minterm DNF of Theorem~\ref{thm:crown}, with the
  resolved/endpoint evaluation table machine-checked;
\item \textbf{the count} --- the explicit bijection
  (\texttt{restrictRep}/\texttt{extendSD} with both inverse laws)
  between self-dual maps of arity $n{+}1$ and free assignments on the
  first-coordinate-\TRU{} representatives;
\item \textbf{calibration} --- \texttt{maj3\_realized} (the explicit
  routing term computes ternary majority; kernel \texttt{decide} over
  the 8-point face) and \texttt{and2\_not\_selfDual};
\item \textbf{one-constant completion} --- \texttt{sdGuard} with
  \texttt{sdGuard\_selfDual}, \texttt{sdGuard\_spec}, and the packaged
  \texttt{one\_constant\_completion}.
\end{itemize}
An independent replay script exhaustively checks the finite claims
(detector tables, the majority term over $R^3$, brute-force crown counts
for $n\le 4$, the $2^n$ miscount check, completion at $n\le 2$, and the
$|R|=6$ counterexample of Section~\ref{sec:sharp}). The frozen output
ships alongside.

\begin{thebibliography}{9}
\bibitem{janov} Yu.~I. Janov, A.~A. Muchnik. \emph{On the existence of
$k$-valued closed classes without a finite basis.} Doklady Akademii Nauk
SSSR, 127:44--46, 1959.
\bibitem{lau} D.~Lau. \emph{Function Algebras on Finite Sets.} Springer,
2006. doi:10.1007/3-540-36023-9.
\bibitem{post} E.~L. Post. \emph{The Two-Valued Iterative Systems of
Mathematical Logic.} Princeton University Press, 1941.
doi:10.1515/9781400882366.
\bibitem{rosenberg} I.~G. Rosenberg. \emph{Complete sets of operations
for finite algebras.} Mathematische Nachrichten, 44(1--6):253--258,
1970. doi:10.1002/mana.19700440120.
\bibitem{wcy2} W.~C. Yang. \emph{The Price of NOT on \textsf{D4}.}
Public edition v1.0, 2026. DOI: 10.5281/zenodo.21800033.
\url{https://github.com/ycmath/price-of-not-on-d4}
\bibitem{wcy1} W.~C. Yang. \emph{Finite-Energy Epistemic Logic with
Conservative Pointed Extension and Negation Geometry.} Public edition
v1.0, 2026. DOI: 10.5281/zenodo.21800031.
\end{thebibliography}

\section*{Authorship \& provenance}

Won Chul Yang, independent researcher. The mathematics of this note is
the author's, developed inside the author's dual-rail carrier programme
around Papers~I--II \cite{wcy1,wcy2}; AI assistance (Anthropic Claude
family) was used for the Lean~4 mechanization, the machine replay, and
the preparation of the manuscript, with the Lean~4 kernel as the
acceptance gate for the verification layer. Remark~\ref{rem:miscount}
records a self-correction of the author's series: an earlier internal
manuscript stated the $2^n$ count; the corrected function count
$2^{2^{n-1}}$ is proved here. Corrections are invited.

Licenses: Apache-2.0 (Lean artifacts and verification scripts),
CC~BY~4.0 (text).

\end{document}
